% !TeX root = ../main.tex
\section{Effect recognition and parameter estimation}
\label{sec:recognition}

This is the most direct line of work: given processed audio, say which effects
are present and with which settings. Almost all of it trains on synthetic
renders, where the true chain is known.

\subsection{Single effects on isolated notes}

Stein et al.\ detected single and cascaded effects on IDMT notes with
hand-crafted spectral, cepstral and harmonic features and support vector
machines, reaching 99.2\% for combinations of six effects on single guitar notes
\reported~\cite{stein2010aes,stein2010dafx}. Such features transfer poorly:
from single notes to polyphonic material, recall fell to chance with one feature
set and to 41--62\% with a bag-of-audio-words model~\cite{schmitt2017recognising}.
The same study noted that delay is trivial to detect on single notes and hard on
phrases, where echoes look like repeated notes. Jürgens et al.\ added parameter
regression with shallow networks (MAE 0.05--0.16 on a unit scale for distortion,
delay and tremolo controls)~\cite{jurgens2020recognizing}. Comunità et al.\
trained a CNN (FxNet) to identify 13 drive plugins (81--91\% accuracy) and a
regressor for gain and tone (MAE 0.02--0.07)~\cite{comunita2021guitar}; most
confusions were between nearly identical circuits such as the TS808 and TS9, a
first sign that the inverse problem is ill-posed at the level of device
identity.

\subsection{Chains and mixes}

Hinrichs et al.\ estimated effect classes and parameters with the guitar still
inside a mix of bass, keys and drums: 97.4\% accuracy on isolated notes and
90.0\% in mixes \reported~\cite{hinrichs2022convolutional}. Guo and McFee
recognised combinations of 13 SoX effects on guitar phrases, with micro and macro
F1 of 0.876 and 0.906 on an unseen dataset~\cite{guo2023automatic}; overdrive vs
distortion, slapback vs feedback delay, and the modulation effects were the usual
confusions. \textbf{Order is harder than presence}: Wada et al.\ reach a set F1
of 0.987 on three-effect chains but a macro F1 of 0.756 on their 16
orders~\cite{wada2025hyperbolic}; Rossi et al.\ an exact-chain accuracy of 0.815
when order and parameters vary~\cite{rossi2025automatic}. Take et al.\ estimated
the chain while recovering the dry signal by peeling effects one at a time;
re-applying the estimated chain to the recovered dry signal improved the wet
reconstruction by only 0.67\dB{} SI-SDR, against 2.24\dB{} from the true dry
signal, because removal and estimation errors add up~\cite{take2024audio}.

\begin{keyidea}[: prediction plus search wins]
Okita and Katayose~\cite{okita2026audio} combined the two approaches of
Section~\ref{sec:search}: a network predicts the dry signal and the effect types,
then CMA-ES or TPE searches the order and parameters of a black-box chain to
maximise the SI-SDR of the re-rendered wet signal. With chorus, distortion and
reverb, the search raised the wet reconstruction from 18.18 to 23.07\dB{} and
lowered MR-STFT from 0.465 to 0.340 \reported. The chorus rate had to be fixed:
its landscape was too rugged. This is the closest published relative of the
app's estimator.
\end{keyidea}

\begin{table}[h]
\centering
\caption{Effect recognition and parameter estimation on guitar, as published.
All training data are synthetic renders.}
\label{tab:recognition}
\small
\begin{tabular}{@{}llll@{}}
\toprule
Work & Input & Task & Result \reported \\
\midrule
Stein~\cite{stein2010dafx} & notes & 6 cascaded effects & 99.2\% acc. \\
Comunità~\cite{comunita2021guitar} & notes & 13 drives, gain, tone & 81--91\%, MAE 0.04 \\
Hinrichs~\cite{hinrichs2022convolutional} & notes in mix & class, params & 90.0\%, MAE $<$0.016 \\
Guo~\cite{guo2023automatic} & phrases & 13 effects & F1 0.876 / 0.906 \\
Rossi~\cite{rossi2025automatic} & DI phrases & 5-effect chain & exact 0.815 \\
Wada~\cite{wada2025hyperbolic} & phrases & ordered chain & set 0.987, order 0.756 \\
Okita~\cite{okita2026audio} & phrases & type, then search & 23.07\dB{} SI-SDR \\
StemFX~\cite{cheng2026stemfx} & separated stems & mixing FX chain & top-1 86.8\% \\
\bottomrule
\end{tabular}
\end{table}

\subsection{Whole processing graphs}

Outside guitar, Lee et al.\ estimate a whole processing graph (nodes,
connections, parameters) from a processed reference~\cite{lee2023blind}; most
errors were wrong orders or processors hidden by a later destructive one. Later
they fit an over-complete differentiable console to dry tracks and the mix and
prune two-thirds of the processors without losing quality~\cite{lee2024searching,lee2025reverse},
using the GRAFX library~\cite{lee2024grafx}. Colonel and Reiss recovered gain,
pan, EQ, delay and reverb of a mix from its raw tracks~\cite{colonel2021reverse}.
These need the dry tracks, but one idea transfers directly: \emph{fit a fixed,
over-complete chain and let the fit switch parts off}. The app's template chain
does exactly that (a pedal whose amount falls below 0.04 is removed).
StemFX is the first system trained to predict effect chains directly on
separated stems without a dry reference: a band-split encoder and a transformer
decoder over 85 effects, trained on pseudo-stems of about 105k
songs~\cite{cheng2026stemfx}. It predicts mixing effects, not guitar rigs, and
guitar sits in its ``other'' stem. LLM2Fx-Tools predicts executable effect chains
as tool calls from dry and processed pairs~\cite{doh2026llm2fxtools}.

\subsection{Sub-problems of a guitar chain}

\begin{description}[style=nextline,leftmargin=1.2cm]
\item[Modulation] Mitcheltree et al.\ extract the LFO of phasers, flangers and
  choruses from processed audio with 1.8--4.9\% error on guitar and 3.1\% on
  unseen real guitar~\cite{mitcheltree2023modulation}. Carson et al.\ fit
  differentiable flangers, choruses and phasers and need a low-frequency
  weighting of the loss to avoid local minima in delay
  time~\cite{carson2026gradient}.
\item[Delay and reverb] Delay time and mix were estimated only on synthetic
  notes~\cite{jurgens2020recognizing,hinrichs2022convolutional}. Blind estimation
  of artificial reverb parameters exists for speech~\cite{lee2022differentiable},
  and blind RT60 estimation is a room-acoustics topic~\cite{gamper2018blind}; we
  found nothing on the reverb of a guitar track inside a mix.
\item[Dynamics] Compressor parameters by hybrid neural regression with
  model-based inversion~\cite{sun2024neural} or gradient-free search in a
  perceptual feature space~\cite{sun2026blackbox}.
\item[The instrument] Pickup position and plucking point can be estimated from
  the spectrum, also through amplifier chains (median errors of 2--7\,mm for the
  pickup)~\cite{mohamad2017pickup,mohamad2017estimating}. The guitar is part of
  the tone that an effects app cannot change, only compensate with EQ.
\end{description}
